\documentclass[11pt]{article}

\usepackage[T1]{fontenc}
\usepackage{lmodern}
\usepackage[a4paper,margin=1in]{geometry}
\usepackage{amsmath,amssymb,amsthm}
\usepackage{booktabs}
\usepackage{enumitem}
\usepackage{microtype}
\usepackage{xurl}
\usepackage{xcolor}
\usepackage[colorlinks=true,allcolors=blue!55!black]{hyperref}

\newtheorem{theorem}{Theorem}
\newtheorem{lemma}{Lemma}
\newtheorem{proposition}{Proposition}
\theoremstyle{remark}
\newtheorem{remark}{Remark}

\newcommand{\Z}{\mathbb{Z}}
\newcommand{\C}{\mathrm{C}}
\newcommand{\SDS}{\operatorname{SDS}}
\newcommand{\corr}{\operatorname{corr}}

\title{Two Small-Order Classification Theorems\\
for Signed Difference Sets}
\author{Nicolas Masselot\\\small Independent researcher, Paris, France\\
\small \texttt{nicolas.masselot@edu.escp.eu}}
\date{September 21, 2026}

\hypersetup{
  pdftitle={Two Small-Order Classification Theorems for Signed Difference Sets},
  pdfauthor={Nicolas Masselot},
  pdfsubject={Signed difference sets and exact finite classification},
  pdfkeywords={signed difference sets, finite abelian groups, quotient enumeration,
    autocorrelation, computer-assisted proof}
}

\begin{document}
\maketitle

\begin{abstract}
A signed difference set assigns coefficients in $\{0,1,-1\}$ to a finite
group so that every nonidentity periodic autocorrelation has the same value.
We settle two clusters of cases that were open in the April 24, 2026 snapshot
of the La Jolla Signed Difference Set Repository.  First, a signed
$(32,20,4)$ difference set exists in an abelian group of order $32$ if and
only if the group is noncyclic.  The six positive cases are given by explicit
coefficient vectors, while the cyclic case is excluded by a complete
$\C_8\to\C_{16}\to\C_{32}$ quotient refinement.  Second, no abelian group
of order $36$ admits a signed $(36,29,4)$ difference set.  This is proved by
three short exhaustive arguments: an empty $\C_{18}$ quotient system excludes
$\C_{36}$ and $\C_2\times\C_{18}$ simultaneously; an empty
$\C_2\times\C_3^2$ quotient system excludes $\C_3\times\C_{12}$; and a
solver-free two-marginal enumeration excludes $\C_6\times\C_6$.  All searches
use exact arithmetic, expose their complete finite search spaces, and are
accompanied by regenerating verification programs.  An external
computational review reproduced the witnesses and the finite computations
supporting these two theorems; its scope is stated separately below.
\end{abstract}

\noindent\textbf{Keywords.}
Signed difference set; finite abelian group; quotient enumeration;
autocorrelation; computer-assisted proof.

\medskip
\noindent\textbf{2020 Mathematics Subject Classification.}
05B10; 05B20; 20K01.

\section{Introduction}

Let $G$ be a finite group of order $v$.  A signed subset of $G$ is a group
ring element
\[
    D=\sum_{g\in G} a_g g, \qquad a_g\in\{0,1,-1\}.
\]
Following Gordon~\cite{gordon}, $D$ is a signed $(v,k,\lambda)$ difference
set when exactly $k$ coefficients are nonzero and
\begin{equation}\label{eq:sds}
    DD^{(-1)}=(k-\lambda)e+\lambda G.
\end{equation}
Equivalently, its periodic autocorrelation is $k$ at the identity and
$\lambda$ at every nonidentity shift.  Signed difference sets contain
ordinary difference sets and circulant weighing matrices as special cases;
further constructions were obtained by He, Chen, and Ge~\cite{hechenge}.

The starting point of this work was the La Jolla Signed Difference Set
Repository~\cite{repository}, frozen at commit
\texttt{e3bf810c5ee6826cf5030f983f6adf23b0ffd20e} on April 24, 2026.
The following two theorems concern ten entries marked \texttt{Open} in that
snapshot.

\begin{theorem}[Order $32$ classification]\label{thm:32}
Let $G$ be an abelian group of order $32$.  A signed $(32,20,4)$ difference
set exists in $G$ if and only if $G$ is noncyclic.  The positive direction is
certified by six explicit constructions; the negative direction is certified
by a complete quotient refinement for $\C_{32}$.
\end{theorem}

\begin{theorem}[Order $36$ classification]\label{thm:36}
No abelian group of order $36$ admits a signed $(36,29,4)$ difference set.
\end{theorem}

The results are existence classifications, not classifications up to
equivalence.  In particular, the six positive files at order $32$ contain one
representative per group; we do not claim that they contain all orbits under
translations, automorphisms, and global sign.

The first version of this note used the repository's cyclic $\C_{36}$
nonexistence and a proof-producing SAT calculation for the final
$\C_6\times\C_6$ case.  Independent review by Ar\'evalo~\cite{arevalo-review}
produced two simplifications now incorporated here: a one-step $\C_{18}$
quotient exclusion for the cyclic case and a direct solver-free exhaustion
for $\C_6\times\C_6$.  The original DRAT certificate is retained as an
independent supplementary check, but neither Theorem~\ref{thm:36} nor its
presentation now depends on SAT.

\section{Arithmetic constraints and quotient projection}

For $h\in G$, write
\[
    \corr_D(h)=\sum_{g\in G}a_g a_{h^{-1}g}.
\]
Equation~\eqref{eq:sds} is equivalent to
$\corr_D(e)=k$ and $\corr_D(h)=\lambda$ for $h\ne e$.
Applying the trivial character gives
\begin{equation}\label{eq:augmentation}
    \left(\sum_{g\in G}a_g\right)^2=k+(v-1)\lambda.
\end{equation}
For every nonprincipal character $\chi$ of an abelian group,
\begin{equation}\label{eq:character}
    \left|\sum_{g\in G}a_g\chi(g)\right|^2=k-\lambda.
\end{equation}
Thus a global sign may normalize the coefficient sum to be nonnegative.
For $(32,20,4)$ the sum is $12$, so a solution has $16$ positive and $4$
negative coefficients.  For $(36,29,4)$ the sum is $13$, so a solution has
$21$ positive and $8$ negative coefficients.

The computational reductions use the following elementary projection.  Its
sum and squared-norm consequences are also instances of the
intersection-number relations in Lemma~5.2 of Gordon~\cite{gordon}.

\begin{lemma}[Quotient projection]\label{lem:projection}
Let $K\trianglelefteq G$ have order $m$, let $Q=G/K$, and define the cell-sum
projection
\[
    B=\sum_{q\in Q}b_q q, \qquad
    b_q=\sum_{g\in q}a_g.
\]
If $D$ is a signed $(v,k,\lambda)$ difference set, then
\begin{equation}\label{eq:quotient}
    BB^{(-1)}=(k-\lambda)e_Q+m\lambda Q.
\end{equation}
Consequently,
\[
    \sum_q b_q=\sum_g a_g,
    \qquad
    \sum_q b_q^2=k+(m-1)\lambda,
\]
and every nonzero quotient shift has correlation $m\lambda$.  Each cell is an
integer in $[-m,m]$.
\end{lemma}

\begin{proof}
The natural map $\Z[G]\to\Z[Q]$ is a ring homomorphism.  Applying it to
Equation~\eqref{eq:sds} sends $e$ to $e_Q$ and sends the formal sum $G$ to
$mQ$, which gives Equation~\eqref{eq:quotient}.  Reading the identity
coefficient gives the sum of squares, and applying the trivial character
gives the total sum.  The cell bound follows because each fiber contains $m$
coefficients in $\{0,1,-1\}$.
\end{proof}

\section{What the exhaustive programs enumerate}\label{sec:method}

We describe the finite searches explicitly because a failed heuristic search
would not prove nonexistence.  For a quotient $Q$ of size $r$ and kernel size
$m$, the basic enumerator performs the following steps.

\begin{enumerate}[label=\textbf{E\arabic*.},leftmargin=3em]
\item Generate every vector $b\in[-m,m]^r\cap\Z^r$ having the required total
sum and squared norm from Lemma~\ref{lem:projection}.  A depth-first traversal
uses necessary feasibility bounds: the remaining cells must be able to
reach the remaining sum and squared norm, and Cauchy--Schwarz requires
$(\text{remaining sum})^2\le
(\text{remaining cells})(\text{remaining norm})$.
\item For each terminal vector, calculate all $r$ quotient correlations
directly and retain it precisely when the identity and every nonidentity
value agree with Equation~\eqref{eq:quotient}.
\item When a retained cell of size $2m$ is refined into two cells of size
$m$, generate every bounded pair whose sum is the parent value.  Repeat until
the final entries lie in $\{0,1,-1\}$, then check the support size and every
full-group correlation.
\end{enumerate}

Every loop is finite and uses exact integer arithmetic.  Translation and
automorphism orbits are used only where their actions are proved to lift to
the next quotient.  Regenerating checkers reconstruct the search spaces
rather than trusting saved candidate counts.  Positive witnesses are checked
by two structurally separate implementations: one uses tuple-valued group
elements and a dictionary operation, while the other uses mixed-radix integer
indices.

\section{The six constructions at order 32}

There are seven abelian groups of order $32$.  The six noncyclic groups all
admit signed $(32,20,4)$ difference sets.  Table~\ref{tab:witnesses} points to
one full coefficient vector for each group.  Every vector contains $16$
entries $+1$, four entries $-1$, and twelve zeros, and both validators return
the correlation vector $(20,4,\ldots,4)$.

\begin{table}[ht]
\centering
\caption{Explicit order-$32$ witnesses in the accompanying data.}
\label{tab:witnesses}
\small
\begin{tabular}{@{}ll@{}}
\toprule
Group & Witness file under \texttt{artifacts/witnesses/}\\
\midrule
$\C_2\times\C_{16}$ & \texttt{sds\_32\_20\_4\_2\_16.json}\\
$\C_4\times\C_8$ & \texttt{sds\_32\_20\_4\_4\_8.json}\\
$\C_2^2\times\C_8$ & \texttt{sds\_32\_20\_4\_2\_2\_8.json}\\
$\C_2\times\C_4^2$ & \texttt{sds\_32\_20\_4\_2\_4\_4.json}\\
$\C_2^3\times\C_4$ & \texttt{sds\_32\_20\_4\_2\_2\_2\_4.json}\\
$\C_2^5$ & \texttt{sds\_32\_20\_4\_2\_2\_2\_2\_2.json}\\
\bottomrule
\end{tabular}
\end{table}

The JSON files record the invariant factors, coordinate convention, full
coefficient vector, positive and negative point sets, and complete
autocorrelation output.  They are also exported as $[P,N]$ coordinate lists
in the format of the La Jolla repository.  Appendix~\ref{app:witnesses}
records the six constructions explicitly.  Since existence is verified by
substitution in Equation~\eqref{eq:sds}, the six files prove the positive
direction of Theorem~\ref{thm:32}.

\section{Excluding the cyclic group of order 32}\label{sec:c32}

Let $D$ be a hypothetical signed $(32,20,4)$ difference set in $\C_{32}$.
Projection first to $\C_8$ uses a kernel of order $4$.  The cells therefore
lie in $[-4,4]$, have sum $12$, squared norm
$20+3\cdot4=32$, and nonzero correlation $16$.  The complete enumeration
from Section~\ref{sec:method} has the counts in Table~\ref{tab:c32}.

\begin{table}[ht]
\centering
\caption{Complete cyclic quotient refinement for $(32,20,4)$.}
\label{tab:c32}
\begin{tabular}{@{}lrr@{}}
\toprule
Stage & Examined & Surviving\\
\midrule
$\C_8$: bounded vectors after sum and norm & $9{,}528$ & $56$\\
Affine orbits on $\C_8$ & $56$ & $5$ orbits\\
$\C_{16}$ refinements of the five representatives & complete & $12$\\
$\C_{32}$ coefficient refinements & $248{,}832$ & $0$\\
\bottomrule
\end{tabular}
\end{table}

The five $\C_8$ orbit sizes are $8,16,8,16,8$.  The affine action is generated
by translations and multiplication by units.  Both operations lift through
$\C_{16}$ to $\C_{32}$ because reduction
$(\Z/32\Z)^\times\to(\Z/8\Z)^\times$ is onto.  For each representative, every
bounded split into two $\C_{16}$ cells is generated and checked.  All twelve
$\C_{16}$ survivors in the normalized parent slices are retained; no second
symmetry quotient is imposed.  Each is then split in every possible way into
two coefficients in $\{0,1,-1\}$, with exact support checked before all $31$
full correlations.  No full vector survives.

As an independent robustness check, Ar\'evalo refined all $56$ $\C_8$
solutions without any orbit reduction.  That calculation produced $144$
$\C_{16}$ survivors and tested $2{,}985{,}984$ final refinements, again with
no solution~\cite{arevalo-review}.  Thus the negative direction does not
depend in practice on the affine reduction.  This completes the proof of
Theorem~\ref{thm:32}.

\section{The order 36 classification}\label{sec:v36}

The abelian groups of order $36$ are
\[
    \C_{36},\qquad \C_2\times\C_{18},\qquad
    \C_3\times\C_{12},\qquad \C_6\times\C_6.
\]
For $(36,29,4)$, Equation~\eqref{eq:augmentation} gives normalized sum $13$.
We exclude the four groups in three computations.

\subsection{One C18 system excludes two groups}

Both $\C_{36}$ and $\C_2\times\C_{18}$ have a quotient $\C_{18}$ with kernel
of order $2$.  A projected vector must therefore lie in $[-2,2]^{18}$, have
sum $13$, squared norm $29+4=33$, and correlation $8$ at every nonzero shift.

Rather than enumerate $5^{18}$ vectors, the verifier constructs the complete
$\C_9$ and $\C_2$ marginal systems first.  The $\C_9$ system has cell bound
$4$, squared norm $41$, and nonzero correlation $16$.  Among $106{,}353$
vectors satisfying sum and norm, its only nine solutions are the nine
translations of
\[
    (-3,2,2,2,2,2,2,2,2).
\]
The $\C_2$ system has precisely the two solutions $(4,9)$ and $(9,4)$.
For every one of the $9\cdot2$ marginal pairs, the program enumerates every
split
\[
    b_i+b_{i+9}=t_i,\qquad b_i,b_{i+9}\in[-2,2],
\]
and retains the required even and odd sums.  There are $19{,}152$
marginal-compatible refinements, of which $7{,}560$ have squared norm $33$.
Direct evaluation of all $18$ cyclic correlations leaves no solution.
Consequently neither $\C_{36}$ nor $\C_2\times\C_{18}$ admits the required
signed difference set.

\subsection{The group C3 x C12}

Project $\C_3\times\C_{12}$ onto $\C_2\times\C_3^2$ with kernel order $2$.
The real marginal is one of $(4,9)$ and $(9,4)$, and the $\C_3^2$ marginal is
a translation of $(-3,2,\ldots,2)$.  Translation in the two quotient factors
acts independently and lifts to $\C_3\times\C_{12}$, so one pair suffices.
The local combined-quotient checker uses this normalization; Ar\'evalo's
external review additionally checked every marginal pair without it.

For the normalized pair, each of the nine $\C_3^2$ cells is split into two
entries in $[-2,2]$, while the row sums are fixed to $(4,9)$.  Exactly $420$
vectors satisfy these marginals and squared norm $33$.  None has correlation
$8$ at all $17$ nonzero shifts of $\C_2\times\C_3^2$.  The quotient system is
empty, so no full signed difference set exists in $\C_3\times\C_{12}$.

\subsection{A direct solver-free search for C6 x C6}

Use the Chinese remainder isomorphism
\[
    \C_6\times\C_6\cong \C_2^2\times\C_3^2.
\]
The complete $\C_3^2$ quotient system consists of the nine translations of
$(-3,2,\ldots,2)$.  The complete $\C_2^2$ quotient system consists of the four
translations of $(7,2,2,2)$.  The direct search examines all $9\cdot4=36$
pairs, so it uses no symmetry normalization.

Fix a marginal pair $(T,R)$.  For each of the nine $\C_3^2$ positions $j$,
generate every pattern $p_j\in\{-1,0,1\}^4$ satisfying
$\sum_z p_j(z)=T_j$.  Split the nine fibers into groups of four and five.
For each half, store the four partial $\C_2^2$ sums; join two half-vectors
exactly when their partial sums add to $R$.  This meet-in-the-middle step is
only an indexing optimization: it generates every full vector with the chosen
marginals exactly once.

Across all $36$ marginal pairs, the joins produce $16{,}964{,}640$ full
vectors.  Each is filtered by the $35$ nonidentity correlation equations in
exact integer arithmetic.  The dot products have magnitude at most $36$, so
the signed $64$-bit batch operations cannot overflow.  Every candidate has
sum $13$; if all $35$ nonidentity correlations were $4$, their sum would force
the identity correlation to be $13^2-35\cdot4=29$.  Thus no separate support
filter is omitted.  No vector survives.  The original SAT encoding
and its checked DRAT trace reach the same conclusion, but are not needed for
this proof.  The four abelian groups have now all been excluded, proving
Theorem~\ref{thm:36}.

\section{Independent replication and verification limits}

Ar\'evalo's review~\cite{arevalo-review} reimplemented the checks from the
definition, independently of the original project's two witness validators,
quotient programs, and SAT encoder.  The present $\C_6\times\C_6$ program
was subsequently adapted from that review, so it is not a third independent
implementation of the direct search.  The review reports the following.

\begin{itemize}
\item All $16$ witnesses from the broader census pass a third checker written
against Gordon's convention.
\item Every count in Sections~\ref{sec:c32} and~\ref{sec:v36} was reproduced.
For the cyclic order-$32$ and the order-$36$ quotient systems, all symmetry
reductions were removed and the same nonexistence decisions were obtained.
\item The direct $\C_6\times\C_6$ search over $16{,}964{,}640$ candidates was
the independent source of the simplified proof now used above.
\item A cell-by-cell comparison of the larger $68$-entry census found $58$
agreements with Ar\'evalo's earlier census and no disagreement; the ten cases
in Theorems~\ref{thm:32} and~\ref{thm:36} are exactly the ten his census had
left open.
\end{itemize}

The review did not audit the CNF encoder or the contents of the large Zenodo
archive.  That limitation no longer bears on Theorem~\ref{thm:36}, because the
direct search replaces the SAT argument.  The review itself used substantial
AI assistance, disclosed by its author; its evidentiary claims are backed by
committed deterministic programs and outputs rather than model responses.

\section{Exact computation and reproducibility}\label{sec:verification}

The accompanying source tree contains the manuscript, witnesses, exact
quotient programs, and machine-readable outputs.  Table~\ref{tab:checks}
connects each part of the two theorems to the finite calculation that checks
it.  The mathematical reductions and completeness arguments are given in
Sections~\ref{sec:method}--\ref{sec:v36}; agreement with a stored count alone
would not establish that a search covers every possible solution.

\begin{table}[ht]
\centering
\small
\caption{Verification targets for the two classification theorems.}
\label{tab:checks}
\begin{tabular}{@{}p{0.30\linewidth}p{0.64\linewidth}@{}}
\toprule
Claim & Exact check and expected result\\
\midrule
Six positive order-$32$ cases & Both validators recompute support $20$ and
all $32$ correlations of each printed construction.\\[3pt]
No solution in $\C_{32}$ & $9{,}528$ bounded $\C_8$ candidates yield $56$
solutions in five affine orbits; $12$ normalized $\C_{16}$ refinements
yield $248{,}832$ final candidates and no solution.\\[3pt]
No solution in $\C_{36}$ or $\C_2\times\C_{18}$ & All $18$ marginal pairs
yield $19{,}152$ refinements; $7{,}560$ have norm $33$, and none satisfies
the $\C_{18}$ correlations.\\[3pt]
No solution in $\C_3\times\C_{12}$ & The normalized quotient system has
$420$ marginal-and-norm candidates and no solution.\\[3pt]
No solution in $\C_6\times\C_6$ & All $36$ marginal pairs yield
$16{,}964{,}640$ full vectors and no solution.\\
\bottomrule
\end{tabular}
\end{table}

From the project root, the following commands rerun the witness audit and
the four finite computations.  The environment requires Python~3; the
direct $\C_6\times\C_6$ calculation additionally uses NumPy~2.0.2.
\begin{verbatim}
.venv/bin/python scripts/audit_final_census.py
.venv/bin/python scripts/verify_v32_remaining_quotients.py
.venv/bin/python scripts/verify_v36_29_4_c18_quotient.py
.venv/bin/python scripts/verify_v36_29_4_combined_quotients.py
.venv/bin/python scripts/verify_v36_29_4_c6x6_direct.py
\end{verbatim}
The first command checks the full census and both witness validators; it
also requires the separate CNF/DRAT files when auditing their stored hashes.
The remaining commands regenerate the cyclic order-$32$ ladder, the empty
$\C_{18}$ system, the combined quotients (including $\C_3\times\C_{12}$),
and the full $\C_6\times\C_6$ search, respectively.  The lightweight package
provides \texttt{verify\_package.py} for checking its file manifest and all
$16$ witnesses without the large trace archive.

The calculations use integer arithmetic, with no floating-point tolerance,
randomized acceptance test, or language-model judgment.  Python integer
operations are unbounded.  The only fixed-width correlation batches use
signed $64$-bit integers on $36$ coefficients in $\{-1,0,1\}$; each product
has magnitude at most one and each sum at most $36$.  Early rejection uses
necessary equations, so a rejected candidate cannot be a solution.  The
saved outputs document the runs, while the source describes how to recreate
their finite search spaces.

These computational checks do not constitute a formal verification of the
Python interpreter, NumPy, or the mathematical reductions.  Repeating an
AI-assisted calculation is also distinct from an independent human proof
audit.  The external review and the independently implemented witness
validators provide the additional checks described above.

\subsection{Relation to the larger census and archived releases}

The two theorems emerged from a census of all $68$ entries with group order at
most $36$ marked \texttt{Open} in the frozen repository snapshot.  That
broader project records $16$ existence results and $52$ nonexistence results.
The ten cases treated here are the principal novelty-supported mathematical
contribution; the other $58$ are independent replications of an exact public
predecessor.  Literature and repository searches found no earlier resolution
of the ten cases, but absence from a search is not a mathematical proof of
novelty.

The public version~1.0 repository is available at~\cite{masselot-repo}.
The original $1.1$~GB collection of CNF formulas and DRAT traces is archived
at Zenodo~\cite{zenodo}.  It remains useful as an independent certificate set
for the larger census, even though the order-$36$ theorem now has a
solver-free proof.  The source package accompanying this submission contains
the additional order-$36$ checks described here; the public version~1.0
release and its DOI should not be mistaken for an archive of this revised
manuscript.

\section{Use of generative AI and author responsibility}

OpenAI Codex was used extensively throughout this project, including to
propose candidate character, quotient, and SAT search strategies; write and
revise search programs, tests, and validators; execute local experiments;
organize evidence and reproducibility materials; and draft and revise the
mathematical exposition.  Its contribution therefore extended to research
exploration and implementation as well as language editing.  The particular
model version used for every historical step is not established by the
project's disclosure record; no single model version is asserted here.

The evidentiary basis is the explicit constructions, complete finite
enumerations, and deterministic checks described in
Section~\ref{sec:verification}.  An AI-generated argument, a failed search,
or an unchecked solver answer is not accepted as evidence.  The two
theorems in this revision use the direct enumerations; the checked DRAT
certificates support additional census entries and a supplementary check
of the $\C_6\times\C_6$ case.  Ar\'evalo's independently developed review
also disclosed substantial AI assistance, with the specific coverage and
limitations stated above.

The human author selected the objective, authorized the computational work,
and reviewed the framing, and remains responsible for the claims and any
submission or public release.  This statement records the documented roles
and checks; it does not assert a separate line-by-line human verification
of every argument or program.  Tan's recent work on circulant weighing
matrices~\cite{tan} provides a related example of reporting exact
computations alongside an explicit AI-use disclosure; the attribution here
describes this project's own workflow.

\paragraph{Acknowledgments.}
The author thanks Daniel M. Gordon for discussing the prior-art search and
mathematical framing, clarifying the repository data format, and encouraging
journal submission.  The author thanks Fabian Ar\'evalo for a detailed
independent reproduction, for the $\C_{18}$ observation, and for the
solver-free $\C_6\times\C_6$ enumeration used in the revised proof.

\paragraph{Supplementary information.}
Online Resource~1 is a reproducibility package containing the exact programs,
witnesses, machine-readable outputs, run instructions, and an integrity
manifest for the computations used in the two theorems.

\section*{Statements and Declarations}

\paragraph{Funding.}
No funding was received to assist with the preparation of this manuscript or
to conduct this study.

\paragraph{Competing interests.}
The author has no relevant financial or non-financial interests to disclose.

\paragraph{Data availability.}
The coefficient vectors, machine-readable outputs, and integrity manifests
supporting the two classification theorems are included in the accompanying
reproducibility package (Online Resource~1).  The version~1.0 CNF and DRAT
archive for the broader census is publicly available from Zenodo at
\url{https://doi.org/10.5281/zenodo.21901581}.  That archive is a historical
release and is not presented as an archive of the present revision.

\paragraph{Code availability.}
The exact Python programs needed to validate the witnesses and regenerate the
finite searches used in the two theorems are included in Online Resource~1.
The public project repository is available at
\url{https://github.com/NicolasMasselot/certified-small-sds-census}; the
submission package records the revision-specific files separately from the
public version~1.0 release.

\paragraph{Author contribution.}
Nicolas Masselot selected the research objective, directed the computational
investigation, reviewed the mathematical framing, and takes responsibility
for the manuscript and its submission.  The substantial roles of generative
AI in research exploration, implementation, experiment execution, evidence
organization, and drafting are described in the preceding section.  Fabian
Ar\'evalo's separately authored review and the components adapted from it are
attributed above and in the cited review artifact.

\appendix
\section{Explicit order-32 constructions}\label{app:witnesses}

Coordinates are written in invariant-factor order.  In each case
$D=\sum_{g\in P}g-\sum_{g\in N}g$.  Direct substitution gives support $20$
and correlation $4$ at every nonidentity shift.

\paragraph{$\C_2\times\C_{16}$.}
\begin{align*}
P={}&\{(0,5),(0,7),(0,10),(0,11),(0,13),(0,15),(1,0),(1,1),\\
&\quad (1,3),(1,4),(1,7),(1,8),(1,9),(1,10),(1,12),(1,15)\},\\
N={}&\{(0,2),(0,3),(1,2),(1,11)\}.
\end{align*}

\paragraph{$\C_4\times\C_8$.}
\begin{align*}
P={}&\{(0,4),(0,5),(0,6),(1,0),(1,1),(1,2),(1,3),(1,5),\\
&\quad (1,6),(2,1),(2,2),(2,4),(3,2),(3,4),(3,6),(3,7)\},\\
N={}&\{(0,0),(2,0),(3,1),(3,5)\}.
\end{align*}

\paragraph{$\C_2^2\times\C_8$.}
\begin{align*}
P={}&\{(0,0,1),(0,0,2),(0,0,3),(0,0,7),(0,1,0),(0,1,3),\\
&\quad (0,1,4),(0,1,6),(0,1,7),(1,0,4),(1,0,5),(1,1,1),\\
&\quad (1,1,3),(1,1,4),(1,1,6),(1,1,7)\},\\
N={}&\{(0,1,1),(1,0,0),(1,0,6),(1,1,0)\}.
\end{align*}

\paragraph{$\C_2\times\C_4^2$.}
\begin{align*}
P={}&\{(0,0,3),(0,1,0),(0,1,1),(0,2,0),(0,2,1),(0,2,2),\\
&\quad (0,3,1),(1,0,0),(1,1,0),(1,1,1),(1,1,3),(1,2,0),\\
&\quad (1,2,3),(1,3,1),(1,3,2),(1,3,3)\},\\
N={}&\{(0,0,0),(0,0,2),(0,1,2),(1,2,1)\}.
\end{align*}

\paragraph{$\C_2^3\times\C_4$.}
\begin{align*}
P={}&\{(0,0,0,0),(0,0,0,2),(0,0,0,3),(0,0,1,1),\\
&\quad (0,0,1,2),(0,0,1,3),(0,1,0,0),(0,1,0,1),\\
&\quad (0,1,1,1),(0,1,1,2),(1,0,0,0),(1,0,0,2),\\
&\quad (1,0,1,2),(1,1,0,1),(1,1,0,2),(1,1,0,3)\},\\
N={}&\{(0,1,0,3),(0,1,1,0),(1,0,0,1),(1,1,1,1)\}.
\end{align*}

\paragraph{$\C_2^5$.}
\begin{align*}
P={}&\{(0,0,1,1,0),(0,0,1,1,1),(0,1,0,0,0),(0,1,0,0,1),\\
&\quad (0,1,0,1,1),(0,1,1,0,1),(1,0,0,0,0),(1,0,0,1,0),\\
&\quad (1,0,0,1,1),(1,0,1,0,0),(1,0,1,1,1),(1,1,0,0,1),\\
&\quad (1,1,0,1,0),(1,1,0,1,1),(1,1,1,0,0),(1,1,1,1,0)\},\\
N={}&\{(0,0,0,1,1),(0,0,1,0,1),(1,0,1,0,1),(1,1,1,0,1)\}.
\end{align*}

\begin{thebibliography}{9}

\bibitem{gordon}
D. M. Gordon,
\newblock Signed difference sets,
\newblock \emph{Designs, Codes and Cryptography} \textbf{91} (2023),
2107--2115.
\newblock \url{https://doi.org/10.1007/s10623-022-01171-8}.

\bibitem{hechenge}
Z. He, T. Chen, and G. Ge,
\newblock New constructions of signed difference sets,
\newblock \emph{Designs, Codes and Cryptography} \textbf{92} (2024),
2323--2340.
\newblock \url{https://doi.org/10.1007/s10623-024-01389-8}.

\bibitem{repository}
D. M. Gordon,
\newblock La Jolla Signed Difference Set Repository, version 1.2,
April 24, 2026.
\newblock \url{https://github.com/dmgordo/signed-difference-sets}.

\bibitem{arevalo-review}
F. Ar\'evalo,
\newblock Bounded review of ``Two Small-Order Classification Theorems for
Signed Difference Sets'',
\newblock independent verification code and report, August 12, 2026,
commit \texttt{05f18c4f9db8d3cb073f78fb486dce7c033a6b69}.
\newblock \url{https://github.com/farev/Matematica/tree/main/conjectures/signed-difference-sets/masselot-review}.

\bibitem{masselot-repo}
N. Masselot,
\newblock Certified Census of Small Signed Difference Sets,
version 1.0 public repository and subsequent project history, 2026.
\newblock \url{https://github.com/NicolasMasselot/certified-small-sds-census}.

\bibitem{zenodo}
N. Masselot,
\newblock Certified Small Signed Difference Sets: CNF and DRAT Proof Archive,
version 1.0, Zenodo, 2026.
\newblock \url{https://doi.org/10.5281/zenodo.21901581}.

\bibitem{tan}
M. M. Tan,
\newblock Character and Multiplier Obstructions for Circulant Weighing
Matrices,
\newblock arXiv:2608.18468v2 [math.CO], September 1, 2026.
\newblock \url{https://arxiv.org/abs/2608.18468v2}.

\end{thebibliography}

\end{document}
